\documentclass[11pt,a4paper]{article}

\usepackage{amsmath, amssymb, amsthm}
\usepackage{geometry}
\usepackage{hyperref}
\usepackage{booktabs}
\usepackage{array}
\usepackage{xcolor}
\usepackage{parskip}
\usepackage{titlesec}
\usepackage{fancyhdr}
\usepackage{microtype}
\usepackage{bm}

\geometry{margin=2.5cm}

\hypersetup{
  colorlinks=true,
  linkcolor=blue!60!black,
  urlcolor=blue!60!black,
}

\pagestyle{fancy}
\fancyhf{}
\rhead{\small CADDack ELBO Reference}
\lhead{\small \nouppercase{\leftmark}}
\cfoot{\thepage}

\titleformat{\section}{\large\bfseries}{{\S}\,\thesection}{1em}{}[\vspace{-4pt}\rule{\linewidth}{0.4pt}]
\titleformat{\subsection}{\normalsize\bfseries}{\thesubsection}{1em}{}

\newcommand{\R}{\mathbb{R}}
\newcommand{\N}{\mathcal{N}}
\newcommand{\code}[1]{\texttt{#1}}
\newcommand{\softplus}{\mathrm{softplus}}
\newcommand{\KL}{\mathrm{KL}}
\newcommand{\NLL}{\mathrm{NLL}}

\title{%
  \textbf{CADDack: The ELBO, End to End}\\[6pt]
  \large What data enters the Bayesian affinity objective, and how it is minimised
}
\author{CADDack Project}
\date{\today}

\begin{document}

\maketitle
\thispagestyle{fancy}

\begin{abstract}
This note traces the training objective of \code{FusionAffinityNet} from raw inputs
to the scalar loss that is minimised. It shows (i) exactly which tensors enter the
model and where, (ii) how they become the predictive parameters $(\mu, s)$, (iii)
the Evidence Lower Bound (ELBO) and the per-batch loss actually computed in
\code{bayes.py:elbo\_loss}, and (iv) the optimisation loop in
\code{train.py:train\_fusion\_from\_complexes}. Every symbol is cross-referenced to
the code. The critical scaling detail --- the KL is divided by the dataset size
$N_{\text{train}}$, not the batch count --- is derived in Section~\ref{sec:scaling}.
\end{abstract}

% ==========================================================================
\section{What data enters, and where}
\label{sec:data}
% ==========================================================================

One training example is a protein--ligand complex
$\mathcal{C} = (\text{ligand SMILES},\, \text{3D geometry},\, y)$, built by
\code{gnn/geometry.py:ComplexExample}. It is turned into \emph{two} independent
graphs plus a scalar label. The affinity model has two deterministic encoders
(``towers'') and one Bayesian head.

\begin{center}
\small
\renewcommand{\arraystretch}{1.35}
\begin{tabular}{@{}p{2.6cm}p{4.2cm}p{2.4cm}p{4.0cm}@{}}
\toprule
\textbf{Symbol} & \textbf{Meaning} & \textbf{Shape} & \textbf{Built by} \\
\midrule
$X_L$ & ligand atom features & $[n_L,\,7]$ & \code{smiles\_to\_graph\_arrays} \\
$\mathcal{E}_L,\,E_L$ & ligand bonds + bond feats & $[2,\,m_L],\,[m_L,\,7]$ & \code{smiles\_to\_graph\_arrays} \\
$z$ & atomic numbers (pocket+ligand) & $[n_G]$ & \code{GeometryRecord} \\
$P$ & 3D coordinates & $[n_G,\,3]$ & \code{GeometryRecord} \\
$b_L,\,b_G$ & batch-assignment vectors & $[n_L],\,[n_G]$ & PyG \code{DataLoader} \\
$y$ & measured affinity (e.g.\ pK\textsubscript{d}) & $[B]$ & \code{geo\_batch.y} \\
\bottomrule
\end{tabular}
\end{center}

\noindent
Here $B$ is the batch size, $n_L$ / $n_G$ the total ligand / geometry atoms in the
batch, and $m_L$ the number of directed ligand edges. The two graphs have
\emph{different} atom counts, so they are carried in two parallel PyG batches and
zipped together (\code{train.py:collate\_pairs}).

\subsection*{Forward pass: from tensors to $(\mu, s)$}

\begin{align}
\mathbf{h}_L &= \textstyle\sum_{\text{pool}} \mathrm{GINE}(X_L, \mathcal{E}_L, E_L, b_L)
              && \in \R^{B\times H} && \code{\_encode\_ligand} \\
\mathbf{h}_G &= \textstyle\sum_{\text{pool}} \mathrm{Geo}(z, P, b_G)
              && \in \R^{B\times H} && \code{GeometryTower} \\
\mathbf{h}   &= [\,\mathbf{h}_L \,\|\, \mathbf{h}_G\,]
              && \in \R^{B\times 2H} && \code{torch.cat} \\
(\mu,\, s)   &= \mathrm{BayesMLP}(\mathbf{h})
              && \mu,s \in \R^{B} && \code{bayes.py:BayesianMLP}
\end{align}

The head returns two vectors per complex: the predicted mean affinity $\mu$ and the
\emph{log-variance} $s = \ln\sigma_{\text{al}}^2$ of the observation noise, clamped to
$[-10, 10]$ for stability (\code{bayes.py:135}). The mean/log-variance heads are
ordinary linear layers; only the hidden layers of the head are Bayesian.

% ==========================================================================
\section{The Bayesian weights}
\label{sec:weights}
% ==========================================================================

Each Bayesian hidden layer places an independent Gaussian over every weight (and
bias) entry $w_{ij}$ (\code{bayes.py:BayesianLinear}):
\begin{equation}
q(w_{ij}) = \N\!\bigl(\mu_{ij},\, \sigma_{ij}^2\bigr),
\qquad
\sigma_{ij} = \softplus(\rho_{ij}) = \ln\!\bigl(1 + e^{\rho_{ij}}\bigr) > 0 .
\end{equation}
The learnable parameters are the pair $(\mu_{ij}, \rho_{ij})$; $\rho$ is used instead
of $\sigma$ directly so $\sigma$ stays strictly positive without a constraint.
A forward pass draws one sample per weight via the reparameterisation trick, which
keeps the draw differentiable in $(\mu,\rho)$:
\begin{equation}
w_{ij} = \mu_{ij} + \sigma_{ij}\,\varepsilon_{ij},
\qquad \varepsilon_{ij}\sim\N(0,1)
\qquad \text{(\code{bayes.py:forward})}.
\end{equation}
The prior is a fixed zero-mean Gaussian $p(w_{ij}) = \N(0,\sigma_{\text{prior}}^2)$
with $\sigma_{\text{prior}} = $ \code{prior\_sigma} (default $1.0$).

% ==========================================================================
\section{The ELBO objective}
\label{sec:elbo}
% ==========================================================================

Variational inference maximises the Evidence Lower Bound over the whole dataset
$\mathcal{D} = \{(X_i, y_i)\}_{i=1}^{N}$, $N \equiv N_{\text{train}}$:
\begin{equation}
\mathcal{L}_{\text{ELBO}}(q)
= \underbrace{\mathbb{E}_{w\sim q}\!\Bigl[\textstyle\sum_{i=1}^{N}\ln p(y_i\mid X_i, w)\Bigr]}_{\text{expected data log-likelihood}}
\;-\; \underbrace{\KL\!\bigl(q(w)\,\|\,p(w)\bigr)}_{\text{complexity penalty}} .
\end{equation}
Maximising $\mathcal{L}_{\text{ELBO}}$ is equivalent to minimising the negative ELBO.
CADDack minimises a per-example--scaled, minibatch estimate of it:
\begin{equation}
\boxed{\;
\mathcal{L}
= \underbrace{\NLL_{\text{batch}}}_{\text{data fit}}
\;+\; \underbrace{\beta \cdot \dfrac{\KL_{\text{total}}}{N_{\text{train}}}}_{\text{regularisation}}
\;}
\qquad \text{(\code{bayes.py:elbo\_loss})}.
\label{eq:loss}
\end{equation}

\subsection{The data-fit term (heteroscedastic Gaussian NLL)}

Assuming $y_i \sim \N\!\bigl(\mu_i,\, e^{s_i}\bigr)$, the negative log-likelihood of a
batch of size $B$, dropping the constant $\tfrac12\ln 2\pi$, is a \emph{mean} over the
batch:
\begin{equation}
\NLL_{\text{batch}}
= \frac{1}{B}\sum_{i=1}^{B}
  \frac{1}{2}\Bigl[\, e^{-s_i}\,(y_i - \mu_i)^2 \;+\; s_i \,\Bigr].
\label{eq:nll}
\end{equation}
The $e^{-s_i}$ factor down-weights noisy examples; the $+\,s_i$ term stops the model
from declaring infinite noise. If \code{no\_aleatoric=True}, this reduces to plain
$\NLL_{\text{batch}} = \frac{1}{B}\sum_i (\mu_i - y_i)^2$ (MSE).

\subsection{The complexity term (closed-form KL)}

Because both $q$ and $p$ are Gaussian, the KL has a closed form per weight; summed
over all Bayesian parameters (\code{bayes.py:kl}):
\begin{equation}
\KL_{\text{total}}
= \sum_{\text{layers}}\;\sum_{i,j}
  \left[\,
    \ln\frac{\sigma_{\text{prior}}}{\sigma_{ij}}
    + \frac{\sigma_{ij}^2 + \mu_{ij}^2}{2\,\sigma_{\text{prior}}^2}
    - \frac{1}{2}
  \right].
\label{eq:kl}
\end{equation}
It is zero exactly when $q = p$ and grows as the posterior moves away from the prior;
$\sigma_{ij}$ is clamped to $\ge 10^{-12}$ before the $\log$ so the term can never
diverge to $-\infty$.

\subsection{Why divide the KL by $N_{\text{train}}$?}
\label{sec:scaling}

The dataset-level objective sums the NLL over all $N$ examples but counts the KL
\emph{once}: $\sum_{i=1}^{N}\NLL_i + \KL$. Dividing the whole thing by $N$ puts it on
a per-example scale,
\begin{equation}
\frac{1}{N}\Bigl(\textstyle\sum_{i=1}^{N}\NLL_i + \KL\Bigr)
= \underbrace{\overline{\NLL}}_{\text{a per-example mean}} + \frac{\KL}{N},
\end{equation}
which is exactly Eq.~\eqref{eq:loss} because $\NLL_{\text{batch}}$ in
Eq.~\eqref{eq:nll} \emph{is} a per-example mean. Hence the KL must be divided by the
number of training examples $N_{\text{train}}$.

\medskip
\noindent\textbf{Common pitfall.} Dividing instead by the number of minibatches
$M = N/B$ (as an earlier version of the code did) leaves the NLL per-example but
scales the KL by $\KL/M = (B)\cdot \KL/N$ --- over-weighting the regulariser by a
factor of the batch size $B$. The posterior then collapses toward the prior
($\mu_{ij}\to 0$), predictions regress to the mean, and epistemic uncertainty is
inflated. CADDack divides by $N_{\text{train}}$ to avoid this.

\subsection{KL warmup ($\beta$ annealing)}

$\beta$ ramps the KL in linearly over the first $\tau$ epochs
(\code{train.py}):
\begin{equation}
\beta(\text{epoch}) = \beta_{\max}\cdot\min\!\Bigl(1,\; \tfrac{\text{epoch}+1}{\tau}\Bigr),
\qquad \beta_{\max}=\code{kl\_weight},\;\; \tau=\code{kl\_warmup}.
\end{equation}
Starting near $\beta\approx 0$ lets the means $\mu_{ij}$ acquire signal before the
regulariser pulls $\sigma_{ij}\to\sigma_{\text{prior}}$; without warmup the network
can get stuck at the prior.

% ==========================================================================
\section{How it is minimised}
\label{sec:minimise}
% ==========================================================================

The loss $\mathcal{L}$ in Eq.~\eqref{eq:loss} is minimised by Adam. Each step uses
\emph{one} weight sample (an unbiased Monte-Carlo estimate of $\mathbb{E}_q[\NLL]$),
while the KL is exact --- no sampling. Gradients flow to $(\mu_{ij}, \rho_{ij})$
through the reparameterised $w_{ij}$ and to the deterministic towers/heads normally.

\begin{quote}
\begin{verbatim}
n_train = number of training complexes        # N
opt     = Adam(model.parameters(), lr)
for epoch in range(epochs):
    beta = kl_weight * min(1, (epoch+1)/kl_warmup)   # KL annealing
    for (lig_batch, geo_batch) in train_loader:      # two-graph batch
        opt.zero_grad()
        mu, log_var = model(lig_batch, geo_batch)    # one weight sample
        y           = geo_batch.y
        kl          = model.kl()                     # closed form, Eq.10
        loss = NLL_batch(mu, log_var, y) + beta*kl/n_train   # Eq.8
        loss.backward()                              # grad wrt mu, rho
        opt.step()
\end{verbatim}
\end{quote}

\noindent This is the training loop of
\code{train.py:train\_fusion\_from\_complexes}. After training, prediction draws
$T$ weight samples and splits the
predictive variance into epistemic ($\mathrm{Var}_t\hat\mu_t$) and aleatoric
($\mathbb{E}_t\,e^{s_t}$) parts (\code{models.py:predict\_with\_uncertainty}); the
biased variance estimator is used so a single sample yields $0$ rather than
\code{NaN}.

\medskip
\noindent\textbf{Gradient of the reparameterised weight} (used implicitly by autograd):
\begin{equation}
\frac{\partial w_{ij}}{\partial \mu_{ij}} = 1,
\qquad
\frac{\partial w_{ij}}{\partial \rho_{ij}}
= \varepsilon_{ij}\,\frac{d}{d\rho}\softplus(\rho_{ij})
= \frac{\varepsilon_{ij}}{1 + e^{-\rho_{ij}}} .
\end{equation}

% ==========================================================================
\section{Symbol / code / shape cross-reference}
% ==========================================================================

\begin{center}
\renewcommand{\arraystretch}{1.3}
\begin{tabular}{@{}llll@{}}
\toprule
\textbf{Math} & \textbf{Code} & \textbf{Shape} & \textbf{Where} \\
\midrule
$\mu$ & \code{mu} & $[B]$ & \code{bayes.py:BayesianMLP.forward} \\
$s=\ln\sigma_{\text{al}}^2$ & \code{log\_var} & $[B]$ & \code{bayes.py:BayesianMLP.forward} \\
$y$ & \code{geo\_batch.y} & $[B]$ & \code{train.py} \\
$\KL_{\text{total}}$ & \code{model.kl()} & scalar & \code{bayes.py:kl} \\
$N_{\text{train}}$ & \code{n\_train} & scalar & \code{train.py:272} \\
$\beta$ & \code{beta} & scalar & \code{train.py:303} \\
$\mathcal{L}$ & \code{loss} & scalar & \code{bayes.py:elbo\_loss} \\
$(\mu_{ij},\rho_{ij})$ & \code{weight\_mu,\ weight\_rho} & layer & \code{bayes.py:BayesianLinear} \\
$\sigma_{\text{prior}}$ & \code{prior\_sigma} & scalar & \code{bayes.py:BayesianLinear} \\
\bottomrule
\end{tabular}
\end{center}

\end{document}
